\section{Complete Experimental Plan and Blank Result Tables}
\label{app:experiments}
\subsection{Reading this section}
All sample counts, caps, and trial counts below are \emph{planned controls}.
They are not completed-result counts. Every numerical measurement cell remains
blank. No plot contains fabricated empirical points. Complete each template
only through a reproducible derivation from immutable raw artifacts, retain
unsuccessful regimes, and list the exact effective method.

\subsection{Datasets, model families, and lanes}
\begin{table*}[t]
\caption{Experimental coverage and native metrics. Target counts are planning
limits, not verified completion counts. Every lane needs pinned revisions.}
\small\begin{tabularx}{\textwidth}{P{2.2cm}P{3cm}P{2.3cm}Y}\toprule
Lane & Model / tasks & Metric & Purpose and admissibility\\\midrule
Primary unified & LLaVA-1.5-7B; MMStar (1,500), TextVQA (3,000), DocVQA (4,500) & Accuracy; consensus; ANLS & Preserve existing primary evidence separately; reuse only exact control matches\\
PrefixKV native & Paper checkpoint; LLaVA-Description; MM-Vet where executable & Corpus PPL; ROUGE-L F1; declared MM-Vet scorer & Reproduce published operating points in the compatible legacy transport before broader claims\\
PrefixKV byte-capped & Same validated checkpoint and tasks within lane & Same native metrics, separate columns & Deterministically rescale profile shape; label byte-matched \texttt{prefixkv\_reimpl} variant\\
VL-Cache native & LLaVA-v1.6-Mistral; COCO Caption; DocVQA & CIDEr; ANLS & Native nominal ratios; local implementation explicitly \texttt{vl\_cache\_reimpl}\\
VL-Cache byte-capped & Same checkpoint and task ledgers as its unified lane & CIDEr; ANLS & Physical-byte caps, not substituted paper ratios\\
LOOK-M & Compatible LLaVA-1.5; MileBench & Accuracy; ROUGE-L; official subtask summaries & Official only when pinned upstream executes; otherwise audited \texttt{lookm\_reimpl}\\
Conditional extensions & Qwen family / multi-image / video / MathVista & Task-specific & Separate authorization, support validation, and preregistration\\
\bottomrule\end{tabularx}
\end{table*}

Different checkpoints cannot be made comparable by renaming their family.
The original legacy LLaVA checkpoint and its HF conversion may differ in
processor, chat conventions, and execution behavior. The model-native lanes
stay separate until checkpoint and preprocessing equivalence are demonstrated.
Unsupported high-retention memory points are explicitly marked hardware
unsupported; do not silently lower their ratio or resolution.

Dataset identity includes repository, configuration, revision, split, and
sample ID. Media hashes and prompt hashes add a second layer of reproducibility.
Questions sharing the same image/page need a declared independence policy.
Calibration, feasibility debugging, and held-out evaluation have distinct
ledgers. Previously inspected samples can support development but cannot be
relabeled untouched. MMStar is not silently replaced by another dataset; any
replacement would require a separately named protocol.

\subsection{Experiment 1: shared-byte quality--memory curves}
For each eligible task/backend/model stratum, run Prompt Top-K, MosaicKV, and
the lane's audited compressed baseline at common physical caps. Hold samples,
generation, preprocessing, attention backend, block granularity, and mandatory
policy fixed. FullKV is run or reused as the quality reference, not as a latency
competitor. At each cap, verify maximal feasible fill and achieved-byte
tolerance before pairing quality or runtime.

Plot each method's measured points as a curve with CIs, without fabricating
missing/unsupported budgets. Connecting observed points visually is not
permission to claim an unmeasured operating point. A strict Pareto comparison
uses observed points with the same task and controls. A confidence-aware plot
should distinguish point-estimate dominance from statistically supported
quality non-inferiority.

\begin{figure}[h]
\centering\begin{tikzpicture}[node distance=6mm]
\node[box,text width=6.2cm] (ledger) {Immutable samples, prompts, generation settings,\\ model revision, backend, and physical-byte cap};
\node[box,text width=1.5cm,below=of ledger.south west,anchor=north west] (a) {Prompt Top-K};
\node[box,text width=1.5cm,below=of ledger] (b) {Audited lane baseline};
\node[box,text width=1.5cm,below=of ledger.south east,anchor=north east] (c) {MosaicKV variant};
\node[box,text width=6.2cm,below=25mm of ledger] (gate) {Check same IDs, realized bytes, controls, failures\\ Keep incompatible rows in separate strata};
\node[box,text width=6.2cm,below=of gate] (stats) {Official task scorer + paired corpus-aware bootstrap\\ Quality--memory frontier and compressed-only latency};
\draw[flow] (ledger.south west)++(0.85,0)--(a.north);
\draw[flow] (ledger)--(b);
\draw[flow] (ledger.south east)++(-0.85,0)--(c.north);
\draw[flow] (a)--(a.south|-gate.north);
\draw[flow] (b)--(gate);
\draw[flow] (c)--(c.south|-gate.north);
\draw[flow] (gate)--(stats);
\end{tikzpicture}

\caption{The direct-comparison pipeline. Every branch receives the same
controls. Shared realized bytes, not matching ratio labels, admit a pair.}
\end{figure}

\begin{table*}[t]
\caption{Shared-byte curve template. Instantiate per model/task/backend and
method; do not average incompatible metrics. Numeric measurement cells blank.}
\small\begin{tabular}{llrrrrrrr}\toprule
Method & Cap & Actual MiB & Padding MiB & Residual MiB & Score & Paired $\Delta$ [CI] & E2E p50 & Paired $n$\\\midrule
Prompt Top-K & 10 MiB & & & & & & & \\
Lane \texttt{\_reimpl} & 10 MiB & & & & & & & \\
MosaicKV exact & 10 MiB & & & & & & & \\
Prompt Top-K & 20 MiB & & & & & & & \\
Lane \texttt{\_reimpl} & 20 MiB & & & & & & & \\
MosaicKV exact & 20 MiB & & & & & & & \\
Prompt Top-K & 40 MiB & & & & & & & \\
Lane \texttt{\_reimpl} & 40 MiB & & & & & & & \\
MosaicKV exact & 40 MiB & & & & & & & \\
Prompt Top-K & 80 MiB & & & & & & & \\
Lane \texttt{\_reimpl} & 80 MiB & & & & & & & \\
MosaicKV exact & 80 MiB & & & & & & & \\
\bottomrule\end{tabular}
\end{table*}

The 30\% target is tested on the quality-constrained frontier, not by comparing
arbitrarily chosen nominal ratios. If the most memory-efficient quality-eligible
competitor requires $B_c$ and MosaicKV requires $B_m$, the relative reduction
is $1-B_m/B_c$. The competitor-selection rule and threshold are frozen before
evaluation. Where MosaicKV requires more memory, show that result plainly.

\subsection{Experiment 2: paper-native reproduction}
\paragraph{PrefixKV.}
First validate one or more published PPL/ROUGE-L operating points using the
paper checkpoint, supplied profile or preregistered disjoint calibration,
teacher-forcing target, and evaluator. A released profile with undisclosed
calibration identities must be labeled upstream-profile provenance unresolved,
not a verified clean local calibration. Record paper retention and actual
active bytes. PPL success does not certify ROUGE-L or common-HF transport.
The anomalously low historical OCR results are diagnostic failures, not evidence
against PrefixKV.

\paragraph{VL-Cache.}
Run the native caption/document metrics with audited H2O, StreamingLLM, and
PyramidKV variants where executable, plus MosaicKV and FullKV quality reference.
Label the local implementation \texttt{vl\_cache\_reimpl}. An implementation
of the paper's scoring formula with a different fused attention kernel can
support algorithmic comparisons but not a claim to reproduce its published
systems speedup. Byte-matched results form a separate named lane.

\paragraph{LOOK-M.}
Use MileBench routing and hierarchy. An official run requires the pinned code
to execute under its supported model/environment; a broken upstream path is
not repaired by relabeling local code as official. Include H2O, SnapKV, or RoCo
only after their provenance and compatibility are audited. Accuracy and ROUGE-L
remain separate, alongside the official hierarchy.

\begin{table*}[t]
\caption{Paper-native reproduction gate template. Paper-reported values are
kept in their own column and never masquerade as local measurements. No values
are populated before source/metric verification.}
\small\begin{tabularx}{\textwidth}{lYYrrrr}\toprule
Method & Provenance & Metric / operating point & Paper reported & Local & Delta & Gate\\\midrule
PrefixKV & official pinned & PPL / preregistered ratio & & & & \\
PrefixKV & official pinned & ROUGE-L / preregistered ratio & & & & \\
PrefixKV variant & \texttt{prefixkv\_reimpl} & Byte-capped native metric & & & & \\
VL-Cache & \texttt{vl\_cache\_reimpl} & CIDEr / native ratio & & & & \\
VL-Cache & \texttt{vl\_cache\_reimpl} & ANLS / native ratio & & & & \\
LOOK-M & official if executable & MileBench declared metric & & & & \\
LOOK-M & \texttt{lookm\_reimpl} & MileBench declared metric & & & & \\
\bottomrule\end{tabularx}
\end{table*}

\subsection{Experiment 3: systems microbenchmarks}
Select a fixed latency subset without inspecting method outcomes. Sweep input
visual-token count, supported image count, output length, batch size, and shared
byte cap one axis at a time or through a declared factorial design. Proposed
output lengths are 32, 64, 128, and 256; batch sizes are 1, 2, and 4 where
supported. These are controlled workload values, not observed outputs.

Use isolated GPU runs where possible, record background processes, and record
fixed clocks only when permission and hardware permit them. Each configuration
has five warmups and at least twenty measured repetitions. Alternate or
randomize method order to reduce thermal and temporal drift while preserving
the paired design. An A100 40 GB and A100 80 GB are distinct hardware strata.
Report request throughput and token throughput with explicit denominators;
batch latency divided by batch size is not individual request latency.

\begin{table*}[t]
\caption{MosaicKV phase and storage template, repeated per task/backend/cap.
Nested phase times must not be double-counted in the end-to-end total.}
\small\begin{tabularx}{\textwidth}{YrrrrY}\toprule
Phase & p50 (s) & p95 (s) & Mean (s) & Peak bytes & State or cost represented\\\midrule
Image/video encoder & & & & & Vision features\\
Projector & & & & & Features mapped to LM space\\
Language prefill & & & & & Original full KV creation\\
Future forecast & & & & & Query capture and optional draft\\
Planner preparation & & & & & Pooling, routing, data movement\\
Sparse graph & & & & & Nodes, edges, top-k construction\\
Utility and selection & & & & & Scores and feasible allocation\\
Tier construction & & & & & Exact and compatible prototypes\\
Cache packing & & & & & Physical layout and validity masks\\
Residual transfer & & & & & Host-to-device restored members\\
Repair re-decode & & & & & Recomputed provisional step\\
Final decode & & & & & Token-by-token generation\\
End-to-end wall time & & & & & Full critical path, not phase sum\\\midrule
Active exact KV bytes & & & & & Dedicated storage observation\\
Active prototype KV bytes & & & & & Only if mechanism active\\
Physical padding bytes & & & & & Not populated source evidence\\
CPU residual bytes & & & & & Include scales/indices separately\\
\bottomrule\end{tabularx}
\end{table*}

The break-even condition is useful as a model, not a result. If planning costs
$T_c$, repair costs $T_r$, and each of $D$ decode steps saves $\Delta t$, then
a speedup requires roughly $D\Delta t>T_c+T_r$, ignoring interactions that must
be measured. Compare this explanation with phase traces. A short-context eager
gain cannot be generalized to paged attention; if the compressed methods are
latency-equivalent under vLLM, report that rather than emphasizing only HF.

\subsection{Experiment 4: component ablations}
Run the same feasible budget for all variants. Disable one component at a time,
keep training-free hyperparameters fixed, and record the effective mechanism.
Use a general visual task, an OCR/document task, ChartQA
\citep{masry2022chartqa}, and Video-MME \citep{fu2024videomme}
only when the model supports it. Keep auxiliary task extensions separate from
the primary three-task study.

\begin{table*}[t]
\caption{Ablation template. Each row requires actual mechanism activity and
identical physical budgets; unsupported rows are not scored as zero.}
\small\begin{tabular}{llrrrrrr}\toprule
ID & Variant & Score & Active bytes & Forecast recall & Graph coverage & Logit KL & E2E (s)\\\midrule
A & Full enabled MosaicKV design & & & & & & \\
B & No future-query forecast & & & & & & \\
C & Prompt-window forecast only & & & & & & \\
D & Draft forecast only & & & & & & \\
E & No evidence graph & & & & & & \\
F & Semantic-only graph & & & & & & \\
G & No value contribution & & & & & & \\
H & Attention-only utility & & & & & & \\
I & No prototype merging & & & & & & \\
J & No residual repair & & & & & & \\
K & Oracle forecast (evaluation only) & & & & & & \\
L & Oracle repair (evaluation only) & & & & & & \\
\bottomrule\end{tabular}
\end{table*}

Use at most three primary ablation figures: forecast recall versus task quality,
graph/value contribution versus fidelity, and prototype/repair quality versus
cost. A graph-coverage increase with unchanged score is reported as such.
Disagreement between attention-output error and logit KL is not hidden behind
one preferred metric. If prototypes fall back to exact selection, their row
does not answer whether merging helped.

\subsection{Experiment 5: residual repair, conditional}
This experiment is admitted only after lossless residual indexing, safe
prototype construction, provisional rollback, budget enforcement, and one-time
re-decode pass their correctness tests. Compare none, entropy, prototype risk,
combined risk, and oracle evaluation-only policies. Sweep thresholds, maximum
restored blocks, residual dtype, and placement on a disjoint calibration set,
then freeze the selected policy.

Report repair precision as beneficial triggered repairs divided by triggered
repairs. Define the unit (event or example) before computing it. Recall requires
an independently measured set of repairable compressed failures; an untested
failure is not automatically repairable. Report recovered FullKV answer,
improved task score, harmful change, and no-effect cases, plus p95 repair time,
restored bytes, and net quality/latency change. Include representative failures
from OCR, chart, ordinary images, and video only if available. Do not select
only visually persuasive successes.

\begin{table}[h]
\caption{Conditional repair template. All empirical cells blank.}
\small\begin{tabularx}{\columnwidth}{Yrrrr}\toprule
Policy & Prec. & Recall & p95 time & $\Delta$ score\\\midrule
None & & & & \\
Entropy & & & & \\
Prototype risk & & & & \\
Combined & & & & \\
Oracle (evaluation only) & & & & \\
\bottomrule\end{tabularx}
\end{table}

\subsection{Experiment 6: model/backend generalization, conditional}
Tune on one model and one disjoint calibration dataset, then freeze the
algorithmic hyperparameters. Evaluate on separately supported LLaVA and Qwen
configurations without retuning. Treat eager, SDPA, FlashAttention-2, and
verified vLLM native paths as separate backend implementations. Report
architecture-specific adaptations such as positional transforms, GQA mapping,
visual token expansion, and block granularity. Unsupported combinations remain
explicit. SGLang is excluded from this thesis's executable scope.

Algorithmic generalization is a quality/memory question under a frozen policy.
Backend numerical parity is a correctness question under a fixed algorithm.
Passing one does not prove the other. Report both in separate tables, with
model-by-task quality loss and backend token/logit deviation.

\subsection{Historical artifact reuse and admissibility}
Reuse requires matching checkpoint revision, processor/chat template, media
preprocessing, exact prompts and sample IDs, generation controls, precision,
backend/kernel, block and mandatory policy, and per-sample realized-byte
boundary. Latency reuse additionally requires compatible hardware, warmups,
repetitions, concurrency, and measurement definitions. A matching method name,
nominal ratio, or sample count is insufficient.

Historical nominal-ratio results remain separate quality--memory/Pareto
diagnostics. They are not relabeled shared-byte evidence. Primary FullKV,
Prompt Top-K, and MosaicKV artifacts remain immutable, and the new reproduction
lanes extend rather than overwrite them. No FullKV latency comparison is
introduced by reuse.

\subsection{Figure plan without fabricated results}
The final manuscript should contain: (1) quality--memory curves by task and
backend, (2) compressed-only paired latency ratios at shared realized bytes,
(3) planning/decode phase breakdowns, and (4) conditional ablation or repair
figures if admissible. Each plot reads measured aggregates and displays failed
or unsupported points separately. This draft contains architectural TikZ
diagrams only; empty result tables establish the reporting structure without
manufacturing a visual trend.
